% 第6章 6.7.2–6.7.3：七点图上的阈值扫描。两个单位权三角形 {1,2,3}、{4,5,6} 由单位权边 3-4 相连，
% 顶点 7 以权 0.6 连到顶点 2；度 d=(2,2.6,3,3,2,2,0.6)，vol(V)=15.2。
% L 的第二特征值 lambda_2=0.3174，第二特征向量（第一分量取正）约为
%   u2=(0.2163,0.2889,0.0751,-0.3037,-0.4449,-0.4449,0.6132)；
% 广义问题 L y = mu D y 的第二特征值 mu_2=0.1814，y2 约为 (0.2581,0.2842,0.1384,-0.2025,-0.3178,-0.3178,0.3472)。
% 两个方向的排序相同：7,2,1,3,4,{5,6}（5 与 6 坐标相同，作为一个整体）。
% (a) 按排序位置画出两种方向，各按最大绝对值缩放到 1：
%   u2: 1, 0.4710, 0.3527, 0.1225, -0.4952, -0.7255, -0.7255
%   y2: 1, 0.8186, 0.7435, 0.3986, -0.5832, -0.9154, -0.9154
% (b) 五个阈值给出的 S：{7},{2,7},{1,2,7},{1,2,3,7},{1,2,3,4,7}
%   RatioCut = 0.7, 1.4, 1.1667, 0.5833, 1.4；NCut = 1.0411, 0.7917, 0.5846, 0.2648, 0.6786。
%   两种目标都在 S={1,2,3,7} 处最小，等于枚举得到的离散最优值；虚线为谱下界。
\begin{tikzpicture}[font=\figurefont\footnotesize]
\pgfplotsset{every axis/.style={width=56mm,height=42mm,axis lines=left,
  x axis line style={draw=FigureStroke,line width=.6pt},y axis line style={draw=FigureStroke,line width=.6pt},
  tick style={draw=FigureStroke},tick label style={font=\figurefont\scriptsize,text=FigureInk},
  label style={font=\figurefont\scriptsize},
  ymajorgrids,grid style={draw=FigureGrid,line width=.4pt},
  title style={font=\figurefont\footnotesize},
  legend style={draw=none,fill=none,font=\figurefont\scriptsize,cells={anchor=west},at={(0.5,-0.42)},anchor=north}}}
\begin{axis}[name=A,xmin=0.5,xmax=7.5,ymin=-1.2,ymax=1.2,xtick={1,...,7},xticklabels={$7$,$2$,$1$,$3$,$4$,$5$,$6$},
  ytick={-1,-0.5,0,0.5,1},xlabel={按坐标排序的顶点},title={(a) 两种谱方向（缩放到最大值$1$）}]
\addplot[draw=FigureMuted,dashed,line width=.6pt,forget plot] coordinates {(0.5,0) (7.5,0)};
\foreach \t in {1.5,2.5,3.5,5.5}{\edef\temp{\noexpand\addplot[draw=FigureMuted,dotted,line width=.6pt,forget plot] coordinates {(\t,-1.2) (\t,1.2)};}\temp}
\addplot[draw=FigureRed,line width=.9pt,forget plot] coordinates {(4.5,-1.2) (4.5,1.2)};
\addplot[draw=FigureBlue,line width=1.1pt,mark=*,mark size=1.5pt,mark options={fill=FigureBlue,draw=FigureBlue}] coordinates {(1,1) (2,0.4710) (3,0.3527) (4,0.1225) (5,-0.4952) (6,-0.7255) (7,-0.7255)};
\addplot[draw=FigureGreen,line width=1.1pt,dashed,mark=square*,mark size=1.4pt,mark options={solid,fill=FigureGreen,draw=FigureGreen}] coordinates {(1,1) (2,0.8186) (3,0.7435) (4,0.3986) (5,-0.5832) (6,-0.9154) (7,-0.9154)};
\legend{$\bm L$的第二特征向量,广义特征向量$\bm L\bm y=\mu\bm D\bm y$}
\end{axis}
\begin{axis}[at={(A.east)},anchor=west,xshift=10mm,xmin=0.5,xmax=5.5,ymin=0,ymax=1.6,xtick={1,...,5},
  xticklabels={{$1$},{$2$},{$3$},{$4$},{$5$}},xlabel={阈值序号（$S$依次加入$7,2,1,3,4$）},ytick={0,0.5,1,1.5},title={(b) 阈值扫描与谱下界},legend columns=2]
\addplot[draw=FigureBlue,line width=1.1pt,mark=*,mark size=1.5pt,mark options={fill=FigureBlue,draw=FigureBlue}] coordinates {(1,0.7) (2,1.4) (3,1.1667) (4,0.5833) (5,1.4)};
\addplot[draw=FigureGreen,line width=1.1pt,mark=square*,mark size=1.4pt,mark options={fill=FigureGreen,draw=FigureGreen}] coordinates {(1,1.0411) (2,0.7917) (3,0.5846) (4,0.2648) (5,0.6786)};
\addplot[draw=FigureBlue,dashed,line width=.7pt] coordinates {(0.5,0.3174) (5.5,0.3174)};
\addplot[draw=FigureGreen,dashed,line width=.7pt] coordinates {(0.5,0.1814) (5.5,0.1814)};
\legend{RatioCut,NCut,$\lambda_2\approx0.317$,$\mu_2\approx0.181$}
\end{axis}
\end{tikzpicture}
